\chapter{Suivi des retours de prélecture}
% Modèle du canevas E221, section 7. À tenir entre la prélecture (27 novembre 2026) et le dépôt final
% (11 décembre 2026). Remarque reçue -> correction -> emplacement -> statut. Dates réelles uniquement.
% NB : les cellules commençant par « [ » sont entourées d'accolades (sinon \midrule lit un argument optionnel).

\begin{table}[htbp]
  \centering
  \caption{Suivi des retours de prélecture (modèle)}
  \label{tab:suivi-prelecture}
  \begin{tabular}{@{}p{0.30\linewidth}p{0.32\linewidth}p{0.30\linewidth}@{}}
    \toprule
    \textbf{Retour reçu} & \textbf{Correction apportée} & \textbf{Emplacement et statut} \\
    \midrule
    {[Remarque et date]} & {[Action réalisée ou justification]} & {[Chapitre ou page, traité ou à discuter]} \\
    \bottomrule
  \end{tabular}
\end{table}

% Le canevas demande aussi, selon le projet : outils de collecte, diagrammes détaillés, preuves de tests,
% guides utiles (installation, comptes de démonstration fictifs, scénario de démonstration, solution
% de secours). Annexes à créer dans annexes/ puis à \input dans main.tex et à appeler dans le texte.
